\subsection{改革、市场与北方危机：能力增长为何不能自动变成安全}

\subsubsection{把财政问题带到县与市}

庆历新政在一零四三至一零四五年间的进退，已经显示出北宋政治的一个尺度：中枢可以提出考课、学校和任用的方案，却未必能把同一套标准稳定送到所有州县。熙宁以后，这个问题被推向更深处。均输要重新安排转运，青苗要进入季节性的借贷，免役要把原先由人身承担的差役换算为钱，市易又试图调节京城及若干市场的货物与信用。它们不是彼此无关的条目，而是一种设想：让财政不只在收税时出现，也能在运输、劳役与货币不足之处调度资源。

这种设想有真实的制度收获。国家取得了更密的簿书、转运和征调渠道，能够把边费、河工与常平问题放在共同的计算中；一部分地方也能借此缓和临时供给的断裂。但收获不等于每项法令都以相同方式抵达。借贷是否自愿、役钱如何核定、地方官以何种考成压力执行，都会改变政策的含义。把《宋史》食货志或《宋会要辑稿》的一条规定直接当成全国经验，和把反对者的奏疏直接当成全国受害实录，同样会遮蔽这些差异。争论之所以尖锐，正因它涉及谁拥有解释账目、确认灾伤和决定征收时点的权力。

\subsubsection{城市不是制度的背景板}

益州交子务的设置，说明朝廷已经面对一个由商人、钱铺、地方财政和跨区域贸易共同造成的信用问题；它不证明全国已经无摩擦地进入“纸币时代”。江南稻作、两浙手工业和开封消费同样应当放在水路、税仓、户籍和季节风险中理解。城市为官僚、军队、工匠、商人、寺院和寄居者提供相遇的场所，也把物价、治安、消防、粮食与水运的压力集中到可见的尺度。

这正是北宋政治经济的得失相连之处。以市场和文书为接口，国家比单靠实物征发更能调动远方资源；它也更依赖地方能够持续报送、保管、折变和转运。四川的一张交子、河北的一段堤防，所连接的不是抽象的繁荣，而是不同地区在国家账目中被看见的方式。灾荒、河决和起事会使这种连接暴露出成本：救济、修防、军需与地方家庭的生计不能在同一张奏牍里自动协调。

\subsubsection{北方的选择与制度的限度}

辽在北方的稳定一度使宋能够把军费、使节和边市放进可预期的框架；女真兴起却改变了这个框架。宋金海上之盟所追求的燕云，并非一块只等宋军接收的领土。那里有辽的南京、长期经营的道路与居民，也有刚刚展开的女真征服。宋廷在联盟、进军和接收之间面对的是时间、情报和后勤的限制。后来把这一选择概括为一人误国，既忽略辽、金双方的行动，也把多年累积的边防财政问题从现场拿走了。

靖康前后的城防、议和与勤王，尤其提醒我们国家能力有不同的速度。开封能汇聚文书、粮食和官员，并不意味着命令能比围城、道路阻断和地方自保更快抵达。北宋留下的制度收益是一个能以财政、考试、法令和运输组织广大区域的朝廷；代价是首都化的军政体系在北方交通被切断时格外脆弱。关于围城中的人数、誓词和责任，宋、金史书及南宋追忆的叙事各有追责与自辩的目的。较可靠的写法，是承认政权在一一二七年失去中枢这一结果，同时把每个戏剧性细节放回它的材料来历。
